\documentclass[10pt,a4paper]{amsart}
\usepackage{amsmath,amssymb,amsthm,graphicx}
\usepackage{verbatim}
\newenvironment{reference}{\comment}{\endcomment}
%\newenvironment{reference}{}{}
\newenvironment{slogan}{\comment}{\endcomment}
\newenvironment{history}{\comment}{\endcomment}
% For commutative diagrams we use Xy-pic
\usepackage[all]{xy}

% We use 2cell for 2-commutative diagrams.
\xyoption{2cell}
% Theorem environments.
%
\theoremstyle{plain}
\newtheorem{theorem}[subsection]{Theorem}
\newtheorem{proposition}[subsection]{Proposition}
\newtheorem{lemma}[subsection]{Lemma}

\theoremstyle{definition}
\newtheorem{definition}[subsection]{Definition}
\newtheorem{example}[subsection]{Example}
\newtheorem{exercise}[subsection]{Exercise}
\newtheorem{situation}[subsection]{Situation}

\theoremstyle{remark}
\newtheorem{remark}[subsection]{Remark}
\newtheorem{remarks}[subsection]{Remarks}

\numberwithin{equation}{subsection}

% Macros
%
\def\lim{\mathop{\mathrm{lim}}\nolimits}
\def\colim{\mathop{\mathrm{colim}}\nolimits}
\def\Spec{\mathop{\mathrm{Spec}}}
\def\Hom{\mathop{\mathrm{Hom}}\nolimits}
\def\Ext{\mathop{\mathrm{Ext}}\nolimits}
\def\SheafHom{\mathop{\mathcal{H}\!\mathit{om}}\nolimits}
\def\SheafExt{\mathop{\mathcal{E}\!\mathit{xt}}\nolimits}
\def\Sch{\mathit{Sch}}
\def\Mor{\mathop{\mathrm{Mor}}\nolimits}
\def\Ob{\mathop{\mathrm{Ob}}\nolimits}
\def\Sh{\mathop{\mathit{Sh}}\nolimits}
\def\NL{\mathop{N\!L}\nolimits}
\def\CH{\mathop{\mathrm{CH}}\nolimits}
\def\proetale{{pro\text{-}\acute{e}tale}}
\def\etale{{\acute{e}tale}}
\def\QCoh{\mathit{QCoh}}
\def\Ker{\mathop{\mathrm{Ker}}}
\def\Im{\mathop{\mathrm{Im}}}
\def\Coker{\mathop{\mathrm{Coker}}}
\def\Coim{\mathop{\mathrm{Coim}}}

% Boxtimes
%
\DeclareMathSymbol{\boxtimes}{\mathbin}{AMSa}{"02}

%
% Macros for moduli stacks/spaces
%
\def\QCohstack{\mathcal{QC}\!\mathit{oh}}
\def\Cohstack{\mathcal{C}\!\mathit{oh}}
\def\Spacesstack{\mathcal{S}\!\mathit{paces}}
\def\Quotfunctor{\mathrm{Quot}}
\def\Hilbfunctor{\mathrm{Hilb}}
\def\Curvesstack{\mathcal{C}\!\mathit{urves}}
\def\Polarizedstack{\mathcal{P}\!\mathit{olarized}}
\def\Complexesstack{\mathcal{C}\!\mathit{omplexes}}
% \Pic is the operator that assigns to X its picard group, usage \Pic(X)
% \Picardstack_{X/B} denotes the Picard stack of X over B
% \Picardfunctor_{X/B} denotes the Picard functor of X over B
\def\Pic{\mathop{\mathrm{Pic}}\nolimits}
\def\Picardstack{\mathcal{P}\!\mathit{ic}}
\def\Picardfunctor{\mathrm{Pic}}
\def\Deformationcategory{\mathcal{D}\!\mathit{ef}}
\usepackage{cleveref}
\crefname{originallemma}{original lemma}{original lemmas}
\Crefname{originallemma}{Original lemma}{Original lemmas}
\makeatletter
\newlabel{lemma-base-change-existence}{{lemma-base-change-existence}{source}{}{sourcepointer-1}{}}
\newlabel{lemma-base-change-existence@cref}{{[originallemma][0][]lemma-base-change-existence}{[1][0][]source}{}{}{}}
\makeatother
\crefname{originallemma}{original lemma}{original lemmas}
\Crefname{originallemma}{Original lemma}{Original lemmas}
\makeatletter
\newlabel{lemma-base-change-quasi-compact}{{lemma-base-change-quasi-compact}{source}{}{sourcepointer-2}{}}
\newlabel{lemma-base-change-quasi-compact@cref}{{[originallemma][0][]lemma-base-change-quasi-compact}{[1][0][]source}{}{}{}}
\makeatother
\crefname{originallemma}{original lemma}{original lemmas}
\Crefname{originallemma}{Original lemma}{Original lemmas}
\makeatletter
\newlabel{lemma-closed-immersion-quasi-compact}{{lemma-closed-immersion-quasi-compact}{source}{}{sourcepointer-3}{}}
\newlabel{lemma-closed-immersion-quasi-compact@cref}{{[originallemma][0][]lemma-closed-immersion-quasi-compact}{[1][0][]source}{}{}{}}
\makeatother
\crefname{originallemma}{original lemma}{original lemmas}
\Crefname{originallemma}{Original lemma}{Original lemmas}
\makeatletter
\newlabel{lemma-closed-immersion-valuative-criteria}{{lemma-closed-immersion-valuative-criteria}{source}{}{sourcepointer-4}{}}
\newlabel{lemma-closed-immersion-valuative-criteria@cref}{{[originallemma][0][]lemma-closed-immersion-valuative-criteria}{[1][0][]source}{}{}{}}
\makeatother
\crefname{originallemma}{original lemma}{original lemmas}
\Crefname{originallemma}{Original lemma}{Original lemmas}
\makeatletter
\newlabel{lemma-composition-existence}{{lemma-composition-existence}{source}{}{sourcepointer-5}{}}
\newlabel{lemma-composition-existence@cref}{{[originallemma][0][]lemma-composition-existence}{[1][0][]source}{}{}{}}
\makeatother
\crefname{originallemma}{original lemma}{original lemmas}
\Crefname{originallemma}{Original lemma}{Original lemmas}
\makeatletter
\newlabel{lemma-composition-quasi-compact}{{lemma-composition-quasi-compact}{source}{}{sourcepointer-6}{}}
\newlabel{lemma-composition-quasi-compact@cref}{{[originallemma][0][]lemma-composition-quasi-compact}{[1][0][]source}{}{}{}}
\makeatother
\crefname{originallemma}{original lemma}{original lemmas}
\Crefname{originallemma}{Original lemma}{Original lemmas}
\makeatletter
\newlabel{lemma-universally-closed-local}{{lemma-universally-closed-local}{source}{}{sourcepointer-7}{}}
\newlabel{lemma-universally-closed-local@cref}{{[originallemma][0][]lemma-universally-closed-local}{[1][0][]source}{}{}{}}
\makeatother
\crefname{originalreference}{original reference}{original references}
\Crefname{originalreference}{Original reference}{Original references}
\makeatletter
\newlabel{morphisms-lemma-reach-points-scheme-theoretic-image}{{morphisms-lemma-reach-points-scheme-theoretic-image}{source}{}{sourcepointer-8}{}}
\newlabel{morphisms-lemma-reach-points-scheme-theoretic-image@cref}{{[originalreference][0][]morphisms-lemma-reach-points-scheme-theoretic-image}{[1][0][]source}{}{}{}}
\makeatother
\crefname{originalreference}{original reference}{original references}
\Crefname{originalreference}{Original reference}{Original references}
\makeatletter
\newlabel{stacks-properties-definition-reduced-induced-stack}{{stacks-properties-definition-reduced-induced-stack}{source}{}{sourcepointer-9}{}}
\newlabel{stacks-properties-definition-reduced-induced-stack@cref}{{[originalreference][0][]stacks-properties-definition-reduced-induced-stack}{[1][0][]source}{}{}{}}
\makeatother
\crefname{originalreference}{original reference}{original references}
\Crefname{originalreference}{Original reference}{Original references}
\makeatletter
\newlabel{stacks-properties-lemma-quasi-compact-stack}{{stacks-properties-lemma-quasi-compact-stack}{source}{}{sourcepointer-10}{}}
\newlabel{stacks-properties-lemma-quasi-compact-stack@cref}{{[originalreference][0][]stacks-properties-lemma-quasi-compact-stack}{[1][0][]source}{}{}{}}
\makeatother
\title{Existence in the valuative criterion for a quasi-compact morphism implies that it is universally closed}
\author{The Stacks Project Authors}
\date{stacks-project:stacks-morphisms.tex:local-snapshot-20261012; GNU FDL 1.2 or later}
\begin{document}
\maketitle
\noindent\textit{Editorial note.} All delivered mathematical text is the author's own, copied from the pinned Stacks Project chapter source sources/stacks-algebraic-stacks-morphisms/pinned-original-source.tex; no statement or proof step was rewritten, reordered or removed.\par\medskip
\noindent\textit{Editorial note.} Preamble provenance: the Stacks Project keeps one single shared preamble for every chapter; this card takes it directly from the staged sources/stacks-algebraic-stacks-morphisms/preamble.tex as three verbatim spans (lines 8-12, 14-18 and 154-217), so no derived file was written and nothing was added or redefined. The wrapper class and the base packages are supplied by the assembler.\par\medskip
\noindent\textit{Editorial note.} Bibliography provenance: neither the counted statement nor its proof contains a citation (this lane rescans the whole span before writing the card), so the card carries no bibliography and no entry was reconstructed.\par\medskip
\noindent\textit{Editorial note.} Selection provenance: this card was opened under the file-based control contract (CONTROL.md), with the mechanical pre-card gate run on its statement and proof interval before the card was written.\par\medskip
\noindent\textit{Allowed standard prerequisites (editorial).}\begin{itemize}
\item The category of dotted arrows and its base change reduction. Used throughout the counted proof (source lines 9706-9730). Reformulates the existence hypothesis as a lifting property for dotted arrows.
\item The valuative criterion for universally closed morphisms through the locality of universal closedness. Invoked inside the counted proof (source lines 9731-9755). Reduces universal closedness to the valuative lifting property.
\item The closed immersion and quasi-compactness statements used to reduce the test diagram. Invoked inside the counted proof (source lines 9756-9765). Allows the diagram to be replaced by one over an integral scheme.
\end{itemize}\medskip
\ifdefined\hypertarget\hypertarget{sourcepointer-1}{}\fi
\noindent\textit{Source reference.} Original lemma lemma-base-change-existence: sources/stacks-algebraic-stacks-morphisms/pinned-original-source.tex, source lines 9383--9383. Navigation pointer only: the counted text refers to the Stacks item lemma-base-change-existence, which is an existing labeled result of this same chapter and is not part of this card. The pointer carries the target's own label as its number and its definition line as its span; its content is not reproduced here.\par\medskip
\ifdefined\hypertarget\hypertarget{sourcepointer-2}{}\fi
\noindent\textit{Source reference.} Original lemma lemma-base-change-quasi-compact: sources/stacks-algebraic-stacks-morphisms/pinned-original-source.tex, source lines 1984--1984. Navigation pointer only: the counted text refers to the Stacks item lemma-base-change-quasi-compact, which is an existing labeled result of this same chapter and is not part of this card. The pointer carries the target's own label as its number and its definition line as its span; its content is not reproduced here.\par\medskip
\ifdefined\hypertarget\hypertarget{sourcepointer-3}{}\fi
\noindent\textit{Source reference.} Original lemma lemma-closed-immersion-quasi-compact: sources/stacks-algebraic-stacks-morphisms/pinned-original-source.tex, source lines 2004--2004. Navigation pointer only: the counted text refers to the Stacks item lemma-closed-immersion-quasi-compact, which is an existing labeled result of this same chapter and is not part of this card. The pointer carries the target's own label as its number and its definition line as its span; its content is not reproduced here.\par\medskip
\ifdefined\hypertarget\hypertarget{sourcepointer-4}{}\fi
\noindent\textit{Source reference.} Original lemma lemma-closed-immersion-valuative-criteria: sources/stacks-algebraic-stacks-morphisms/pinned-original-source.tex, source lines 9427--9427. Navigation pointer only: the counted text refers to the Stacks item lemma-closed-immersion-valuative-criteria, which is an existing labeled result of this same chapter and is not part of this card. The pointer carries the target's own label as its number and its definition line as its span; its content is not reproduced here.\par\medskip
\ifdefined\hypertarget\hypertarget{sourcepointer-5}{}\fi
\noindent\textit{Source reference.} Original lemma lemma-composition-existence: sources/stacks-algebraic-stacks-morphisms/pinned-original-source.tex, source lines 9396--9396. Navigation pointer only: the counted text refers to the Stacks item lemma-composition-existence, which is an existing labeled result of this same chapter and is not part of this card. The pointer carries the target's own label as its number and its definition line as its span; its content is not reproduced here.\par\medskip
\ifdefined\hypertarget\hypertarget{sourcepointer-6}{}\fi
\noindent\textit{Source reference.} Original lemma lemma-composition-quasi-compact: sources/stacks-algebraic-stacks-morphisms/pinned-original-source.tex, source lines 1994--1994. Navigation pointer only: the counted text refers to the Stacks item lemma-composition-quasi-compact, which is an existing labeled result of this same chapter and is not part of this card. The pointer carries the target's own label as its number and its definition line as its span; its content is not reproduced here.\par\medskip
\ifdefined\hypertarget\hypertarget{sourcepointer-7}{}\fi
\noindent\textit{Source reference.} Original lemma lemma-universally-closed-local: sources/stacks-algebraic-stacks-morphisms/pinned-original-source.tex, source lines 2746--2746. Navigation pointer only: the counted text refers to the Stacks item lemma-universally-closed-local, which is an existing labeled result of this same chapter and is not part of this card. The pointer carries the target's own label as its number and its definition line as its span; its content is not reproduced here.\par\medskip
\ifdefined\hypertarget\hypertarget{sourcepointer-8}{}\fi
\noindent\textit{Source reference.} Original reference morphisms-lemma-reach-points-scheme-theoretic-image: sources/stacks-algebraic-stacks-morphisms/pinned-original-source.tex, source lines 2141--2141. Navigation pointer only: the counted text refers to the Stacks item morphisms-lemma-reach-points-scheme-theoretic-image, which belongs to another chapter of the project and is not part of this card. The pointer carries the target's own label as its number because this chapter's staged snapshot does not contain the other chapter's tag table; its content is not reproduced here.\par\medskip
\ifdefined\hypertarget\hypertarget{sourcepointer-9}{}\fi
\noindent\textit{Source reference.} Original reference stacks-properties-definition-reduced-induced-stack: sources/stacks-algebraic-stacks-morphisms/pinned-original-source.tex, source lines 6968--6968. Navigation pointer only: the counted text refers to the Stacks item stacks-properties-definition-reduced-induced-stack, which belongs to another chapter of the project and is not part of this card. The pointer carries the target's own label as its number because this chapter's staged snapshot does not contain the other chapter's tag table; its content is not reproduced here.\par\medskip
\ifdefined\hypertarget\hypertarget{sourcepointer-10}{}\fi
\noindent\textit{Source reference.} Original reference stacks-properties-lemma-quasi-compact-stack: sources/stacks-algebraic-stacks-morphisms/pinned-original-source.tex, source lines 1935--1935. Navigation pointer only: the counted text refers to the Stacks item stacks-properties-lemma-quasi-compact-stack, which belongs to another chapter of the project and is not part of this card. The pointer carries the target's own label as its number because this chapter's staged snapshot does not contain the other chapter's tag table; its content is not reproduced here.\par\medskip
\begin{lemma}
\label{lemma-quasi-compact-existence-universally-closed}
Let $f : \mathcal{X} \to \mathcal{Y}$ be a morphism of algebraic stacks.
Assume
\begin{enumerate}
\item $f$ is quasi-compact, and
\item $f$ satisfies the existence part of the valuative criterion.
\end{enumerate}
Then $f$ is universally closed.
\end{lemma}

\begin{proof}
By Lemmas \ref{lemma-base-change-quasi-compact}
and \ref{lemma-base-change-existence}
properties (1) and (2) are preserved under
any base change. By Lemma \ref{lemma-universally-closed-local}
we only have to show that $|T \times_\mathcal{Y} \mathcal{X}| \to |T|$
is closed, whenever $T$ is an affine scheme mapping into $\mathcal{Y}$.
Hence it suffices to show: if $f : \mathcal{X} \to Y$ is a
quasi-compact morphism from an algebraic stack to an affine scheme
satisfying the existence part of the valuative criterion,
then $|f|$ is closed. Let $T \subset |\mathcal{X}|$ be a closed subset.
We have to show that $f(T)$ is closed to finish the proof.

\medskip\noindent
Let $\mathcal{Z} \subset \mathcal{X}$ be the reduced induced
algebraic stack structure on $T$ (Properties of Stacks,
Definition \ref{stacks-properties-definition-reduced-induced-stack}).
Then $i : \mathcal{Z} \to \mathcal{X}$ is a closed immersion
and we have to show that the image of $|\mathcal{Z}| \to |Y|$
is closed. Since closed immersions are quasi-compact
(Lemma \ref{lemma-closed-immersion-quasi-compact})
and satisfies the existence part of the valuative criterion
(Lemma \ref{lemma-closed-immersion-valuative-criteria})
and since compositions of quasi-compact morphisms are quasi-compact
(Lemma \ref{lemma-composition-quasi-compact})
and since compositions preserve the property of satisfying
the existence part of the valuative criterion
(Lemma \ref{lemma-composition-existence})
we conclude that it suffices to show: if $f : \mathcal{X} \to Y$
is a quasi-compact morphism from an algebraic stack to an affine scheme
satisfying the existence part of the valuative criterion,
then $|f|(|\mathcal{X}|)$ is closed.

\medskip\noindent
Since $\mathcal{X}$ is quasi-compact (being quasi-compact over the
affine $Y$), we can choose an affine scheme $U$ and a surjective
smooth morphism $U \to \mathcal{X}$
(Properties of Stacks, Lemma
\ref{stacks-properties-lemma-quasi-compact-stack}).
Suppose that $y \in Y$ is in the closure of the image of $U \to Y$
(in other words, in the closure of the image of $|f|$).
Then by
Morphisms, Lemma \ref{morphisms-lemma-reach-points-scheme-theoretic-image}
we can find a valuation ring $A$ with fraction field $K$
and a commutative diagram
$$
\xymatrix{
\Spec(K) \ar[r] \ar[d] & U \ar[d] \\
\Spec(A) \ar[r] & Y
}
$$
such that the closed point of $\Spec(A)$ maps to $y$. By assumption
we get an extension $K'/K$ and a valuation ring $A' \subset K'$
dominating $A$ and the dotted arrow in the following diagram
$$
\xymatrix{
\Spec(K') \ar[r] \ar[d] &
\Spec(K) \ar[r] \ar[d] &
U \ar[d] \ar[r] &
\mathcal{X} \ar[d]^f \\
\Spec(A') \ar[r] \ar@{..>}[rrru] &
\Spec(A) \ar[r] &
Y  \ar@{=}[r] &
Y
}
$$
Thus $y$ is in the image of $|f|$ and we win.
\end{proof}

\end{document}
